%%═══════════════════════════════════════════════════════════════════
%% Lecture 21 — Monte Carlo Methods
%% 77315 — Computational Physics
%%═══════════════════════════════════════════════════════════════════

\section{Lecture 21 — Monte Carlo Methods}\label{sec:lec21}
\hrule\vspace{1.5em}

%%──────────────────────────────────────────────────────────────────
\subsection*{Overview}
Monte Carlo (MC) methods estimate integrals and other quantities
via \emph{random sampling}.  The error decreases as $N^{-1/2}$
regardless of dimension $d$, making MC the method of choice for
$d \ge 4$ compared to deterministic quadrature (which scales as
$N^{-k/d}$).  Variance reduction via \emph{importance sampling}
and the \emph{rejection (von Neumann)} method are introduced.

%%──────────────────────────────────────────────────────────────────
\subsection{MC Integration as a Statistical Average}\label{subsec:lec21:basic}

We wish to evaluate
\begin{equation}
  I = \int_D f(\mathbf{x})\,d^d x.
  \label{eq:lec21:integral}
\end{equation}
Draw $N$ uniform samples $\mathbf{x}_1,\ldots,\mathbf{x}_N$ from
domain $D$ of volume $V_D$:
\begin{equation}
  I \approx \hat{I} = \frac{V_D}{N}\sum_{i=1}^{N}f(\mathbf{x}_i).
  \label{eq:lec21:estimate}
\end{equation}
$\hat{I}$ is an unbiased estimator with standard error
\begin{equation}
  \sigma_{\hat{I}} = V_D\,\frac{\sigma_f}{\sqrt{N}},
  \quad
  \sigma_f^2 = \langle f^2\rangle - \langle f\rangle^2.
  \label{eq:lec21:error}
\end{equation}
The error scales as $N^{-1/2}$ independent of $d$.

%%──────────────────────────────────────────────────────────────────
\subsection{Comparison with Quadrature}\label{subsec:lec21:comparison}

For a $d$-dimensional integral discretised with $n$ points per
dimension, deterministic quadrature of order $k$ achieves error
$\bigO{n^{-k}} = \bigO{N^{-k/d}}$ where $N=n^d$ is the total number
of points.  MC achieves $\bigO{N^{-1/2}}$ always.

\begin{center}
\begin{tabular}{lcc}
\toprule
Method & Error & Break-even $d$ \\
\midrule
Trapezoid ($k=2$) & $\bigO{N^{-2/d}}$ & $d>4$ \\
MC                & $\bigO{N^{-1/2}}$ & always $N^{-1/2}$ \\
\bottomrule
\end{tabular}
\end{center}

MC is better than trapezoid for $d > 4$.

%%──────────────────────────────────────────────────────────────────
\subsection{Importance Sampling}\label{subsec:lec21:importance}

Choose a \emph{weight function} (probability density) $w(\mathbf{x})>0$
that mimics the shape of $f$.  Rewrite:
\begin{equation}
  I = \int_D \frac{f(\mathbf{x})}{w(\mathbf{x})}\,w(\mathbf{x})\,d^d x
    = \left\langle \frac{f}{w} \right\rangle_w.
  \label{eq:lec21:importance}
\end{equation}
Sample $\mathbf{x}_i$ from $w$ and estimate
$I \approx N^{-1}\sum_i f(\mathbf{x}_i)/w(\mathbf{x}_i)$.

The variance of $f/w$ is minimised when $w \propto |f|$, reducing
the error by orders of magnitude over uniform sampling.

\nt{Importance sampling requires an efficient method to draw samples
from $w$.  For simple $w$ (Gaussian, exponential) the inverse-CDF
transform is used.  For general $w$, use the rejection method
or Metropolis algorithm.}

%%──────────────────────────────────────────────────────────────────
\subsection{Rejection Method (von Neumann)}\label{subsec:lec21:rejection}

To sample from a target $p(x) \propto f(x)$ on $[a,b]$:

\begin{enumerate}
  \item Find $M = \max_{x\in[a,b]} f(x)$.
  \item Draw $x \sim \text{Uniform}(a,b)$ and $y \sim \text{Uniform}(0,M)$.
  \item \textbf{Accept} $x$ if $y \le f(x)$; otherwise \textbf{reject}.
  \item Accepted samples are distributed according to $p(x)$.
\end{enumerate}

The acceptance rate is $A = \int f(x)\, dx / [M(b-a)]$.  For
multidimensional sampling, enclose $f$ in a box of known volume.

\nt{The rejection method is inefficient when $f$ is sharply peaked
($A \ll 1$).  Importance sampling or Markov-chain methods
(Metropolis) are preferred for such distributions.}
